\section{Deployment constraints and decision-gated implementation}\label{sec:deployment}

A promising equation is not a construction permit. Deployment is a sequence of tests in which a failure can force a smaller design, another site or cancellation. The roadmap below is a dependency chain, not a prediction of when Singapore will deploy nuclear power.

The preferred mature numerical case is not a deployment claim. Implementation follows the E7 falsification pathway:
\begin{enumerate}
\item Current research and evidence-gap closure.
\item Component and nuclear/process-interface qualification.
\item First-of-a-kind (FOAK) process-heat demonstration.
\item Two-train 353.6 MWth process-heat demonstration.
\item Replication and early-commercial (BOAK) cost validation.
\item Mature 10-OAK economic decision gate.
\item Singapore site, licensing and contract closure.
\item Construction and commissioning.
\item Commercial operation and measured verification.
\end{enumerate}
This sequence is a logical set of evidence gates, not a calendar schedule; failure at a gate requires redesign, relocation, delay or termination as appropriate.

\subsection{Binding implementation gates}
Before a real Singapore project can proceed, the following must close:
\begin{itemize}
\item \textbf{national/regulatory readiness:} an applicable power-reactor policy, competent regulator, licensing/legal framework, safeguards/security, waste/spent-fuel, EPR, liability and workforce arrangements;
\item \textbf{technology qualification:} 600 MWth coupled transients, high-temperature IHX/secondary-He/reformer qualification, tritium, controls/isolation, materials and inspection;
\item \textbf{two-train proof:} physical delivery of 353.6 MWth within the source process-heat branch, including one- and two-train trips, turndown, restart, CCS interaction and H$_2$ quality;
\item \textbf{integrated safety:} PRA, mechanistic source term, IHX leak/rupture, tritium, chemical HAZOP/QRA, propagation/common cause, UHS/SBO, dispersion/dose, EPR/EPZ and security;
\item \textbf{site/Jurong:} parcel, geotechnical/coastal/external hazards, QRA/PRA-derived separation, protected area, emergency access, cooling and environmental baseline;
\item \textbf{infrastructure/market:} credible H$_2$ offtake for approximately 195,892 t/y, firm 68 MMSCFD NG supply, water/cooling, and contracted permanent storage for at least approximately 1.085 MtCO$_2$/y;
\item \textbf{economic maturation:} observed replication cost, construction, operations and maintenance (O\&M), and integrated availability must support a recalculated cost below S\$100/tCO$_2$e, including known site, licensing, security and CCS costs;
\item \textbf{commercial/FID:} full project CAPEX/TCI, financing/IDC, insurance/liability, decommissioning/waste, EPC/owner costs and binding contracts.
\end{itemize}

Jurong failure does not automatically falsify the technology: another Singapore context may be evaluated. Conversely, if no Singapore context can satisfy siting, safety, security and emergency requirements, Singapore deployment is stopped.

\subsection{STOP/REDESIGN logic}
The project does not advance automatically. It is redesigned, relocated, reduced in scale, delayed or stopped if high-temperature coupling cannot be qualified; the two-train dynamics fail; integrated availability invalidates the model; observed mature economics remain at or above S\$100/tCO$_2$e; nuclear/chemical separation is infeasible; source-term/dose or external hazards are incompatible with a site; no credible UHS exists; H$_2$ offtake or the CCS chain cannot be contracted; or the future Singapore regulator rejects the design/site.

Only measured commissioning and commercial operation can convert the projected BOAK/10-OAK performance into demonstrated performance.

\section{Limitations}
The final process balance is anchored to INL Aspen studies rather than a project-owned detailed equipment simulation. The repository independently reconstructs annualisation, lifecycle extensions, helium-flow screening, deployment economics and CN4252 classification, but it does not reproduce INL's proprietary Aspen stream files. Reactor/IHX integration is screening-level: the 176.8 MWth final duty exceeds the 170 MWth reference GTHTR300C IHX, so the doubled-IHX cost sensitivity is used conservatively without claiming a completed exchanger design.

Upstream-gas, nuclear-lifecycle and CCS-transport factors are screening proxies rather than Singapore-specific measured inventories. Historical-cost escalation uses transparent broad deflators rather than nuclear-specific construction indices. Electricity value and T\&S remain scenarios, not guaranteed Singapore contracts. The mature GTHTR300C cost basis is a design-study estimate rather than an observed commercial plant cost. Nuclear licensing, siting, high-temperature process qualification and cross-border CCS remain future deployment conditions.

Accordingly, any CN4252 pass is a conditional model result under the declared literature/design-study assumptions. It does not establish bankable project feasibility or a technology preference over alternatives whose Singapore economics are not matched in this study.

\section{Nuclear engineering feasibility and safety}
\label{sec:nuclear-safety}

\subsection{Evidence maturity}
JAEA's 30 MWth HTTR is a constructed helium-cooled, graphite-moderated test reactor that achieved a 950$^\circ$C outlet temperature at full power~\cite{jaeahttr2026ops}. The selected 600 MWth GTHTR300C-class configuration is instead a design/economic architecture~\cite{nishihara2007potential}, while INL TEV-961 is a process model showing that 925$^\circ$C reactor outlet and 900$^\circ$C delivered heat can supply the 176.8 MWth reforming duty~\cite{inltev961}. The project therefore combines demonstrated temperature precedent, designed reactor architecture and modelled process integration; it does not claim a demonstrated 600 MWth nuclear-SMR plant.

\subsection{TRISO and passive/inherent characteristics}
\textbf{Plain-language takeaway:} each TRISO particle is a tiny fuel kernel wrapped in several engineered barrier layers; those layers strongly retain fission products, but they are not indestructible. TRISO fuel is a layered pressure vessel at particle scale. Around the fuel kernel, a porous carbon buffer accommodates fission gas and kernel swelling; the inner pyrolytic-carbon (IPyC, a dense carbon coating) layer provides structural support and protects the silicon-carbide deposition interface; dense silicon carbide (SiC, a ceramic barrier) is the principal pressure-bearing and metallic-fission-product barrier; and the outer pyrolytic-carbon (OPyC) layer provides additional mechanical/environmental protection~\cite{sawa2001httrfuel,sawaueta2004}. INL's AGR qualification programme adds modern UCO irradiation and post-irradiation safety evidence: 1600--1800$^\circ$C tests measured fission-product release and identified the consequences of SiC/coating failures~\cite{demkowicz2015facs,morris2016agr1,inlagrrelease}. Intact SiC showed excellent cesium retention in these tests, but measurable failure/release mechanisms remain; TRISO therefore reduces source-term potential (the amount of radioactive material available for release) rather than making radionuclide release impossible.

HTGR safety also uses reactor-scale physics. JAEA reports HTTR average power density of 2.5 MW/m$^3$ and a large graphite inventory~\cite{jaeahttr2026ops}. In the 9 MW HTTR loss-of-forced-cooling test, primary circulators were stopped without control-rod insertion; test/analysis evidence shows power falling rapidly toward decay heat through negative temperature-reactivity feedback while temperatures changed slowly because of graphite thermal inertia~\cite{takamatsu2011lofc}. These features reduce transient severity but do not remove the need for shutdown systems, decay-heat analysis, confinement, emergency planning or accident assessment.

\subsection{Accident mechanisms retained}
\textbf{Plain-language takeaway:} favourable HTGR physics does not remove accident analysis; loss of helium pressure, unwanted air or steam entry, loss of heat removal and radioactive-material transport still have to be evaluated. NRC HTGR/TRISO PIRTs explicitly retain depressurisation accidents with air ingress and water ingress, while the NGNP process-heat PIRT identifies intermediate-loop blowdown/loss of heat sink, leakage into the primary system and reactive-gas ingress as coupled-system safety phenomena~\cite{nrctriso6844,nrcngnppirt2010}. Depressurisation and subsequent air ingress can expose hot graphite to oxygen, so graphite oxidation and its effects on structural integrity remain accident mechanisms; fission-product transport during normal and accident conditions is itself a dedicated NRC PIRT topic~\cite{nrcngnpfpt2008}. Water/steam ingress, loss of heat sink and primary leakage likewise require design-specific analysis. The project therefore does not describe HTGRs as risk free or claim that severe fuel damage is impossible under every condition.

\subsection{IHX: heat-transfer equipment and safety boundary}
\textbf{Plain-language takeaway:} the IHX is both a heater and a physical separator; its metal boundary must survive high temperature for long periods without allowing the reactor and chemical loops to mix. JAEA's constructed HTTR IHX is a 10 MW helium--helium exchanger whose Hastelloy XR tubes form the pressure boundary between primary and secondary helium. Development work addressed creep collapse and creep-fatigue (slow high-temperature deformation combined with repeated thermal/mechanical cycling), seismic behaviour, bundle thermal hydraulics and in-service inspection~\cite{hamamoto2006ihx}. Structural analysis reports primary helium near 950$^\circ$C and secondary helium up to about 905$^\circ$C, with high-temperature creep explicitly governing design~\cite{httrihx1995}.

IAEA coupling guidance identifies the intermediate circuit as a means to separate the nuclear island from the chemical plant, limit radioactive contamination and reduce propagation of chemical disturbances into the reactor~\cite{iaeanuclearhydrogen2013}. This supports the project's primary-helium $\rightarrow$ IHX $\rightarrow$ secondary-helium $\rightarrow$ reformer architecture. It does not remove scale-up work: a constructed 10 MW IHX and a designed 370 MWth source heat branch are not a completed 176.8 MWth project exchanger.

\subsection{Coupled nuclear--chemical transients}
\textbf{Plain-language takeaway:} if one plant suddenly changes state, the other must not be driven into an unsafe condition; ``propagation'' here means one plant's disturbance causing a consequential event in the other. A chemical-plant trip or rapid heat-demand change alters secondary-helium return conditions. IAEA/JAEA work identifies return-helium control and prevention of reactor trips as explicit integration-development needs~\cite{iaeacoupling2009}. Conversely, reactor trip or loss of nuclear heat requires the chemical plant to reach a safe state with hot, pressurised methane/steam/hydrogen inventories. Independent trips, isolation, secondary-loop thermal management and process depressurisation therefore require future transient design; this screening study does not size those systems.

\subsection{Chemical/process hazards and co-location}
The hydrogen plant introduces combustible hydrogen and methane, toxic carbon monoxide, hot high-pressure syngas/steam, reformer fire/explosion hazards, solvent/capture equipment and compressed CO$_2$. A future process hazard analysis must evaluate leakage, ignition, explosion overpressure, toxic/asphyxiant releases, escalation and loss of utilities. The NRC NGNP process-heat/hydrogen PIRT emphasises chemical inventories, layout and separation distance as major controls against propagation to the nuclear plant~\cite{nrcngnppirt2010}; IAEA work likewise identifies combustible leakage/explosion as an HTGR-hydrogen integration issue~\cite{iaeacoupling2009}.

\subsection{Safety-feasibility conclusion}
At screening level the design has a credible physical safety architecture: TRISO fuel, graphite/helium HTGR characteristics, an intermediate helium boundary and demonstrated high-temperature IHX precedent. \textbf{Safety feasibility remains conditional}: site-specific accident consequences, chemical-to-nuclear propagation limits, separation distances, detailed IHX integrity, emergency planning and regulatory acceptance are not demonstrated here.
\section{Singapore implementation implications}
The corrected model translates the candidate into annual H2, fresh-NG energy, captured/residual CO2, cross-border storage throughput, nuclear process heat and helium-loop scale. Required thermal service is expressed against the source-backed 170 MW JAEA IHX benchmark rather than converted into a fabricated commercial reactor-module count.

Singapore deployment remains conditional~\cite{singaporenuclear,singaporeccs}. Nuclear power has not been assumed to be deployed, and cross-border CO2 transport/storage is treated as infrastructure under development rather than an existing domestic storage service. An implementation roadmap would therefore require, at minimum, nuclear regulatory/readiness development, site and safety assessment, a demonstrated process-heat coupling architecture, and contracted cross-border CO2 transport/storage.

